\mysection{ネットワークアドレス計算の例}

ご存知のように、TCP/IPアドレス（IPv4）は$0 \ldots 255$の範囲の4つの数字、つまり4バイトで構成されています。

4バイトは32ビット変数に簡単に収まるため、IPv4ホストアドレス、ネットワークマスク、またはネットワークアドレスはすべて32ビット整数にできます。

ユーザーの観点から見ると、ネットワークマスクは4つの数字として定義され、255.255.255.0などのように形式化されますが、ネットワークエンジニア（システム管理者）は\q{/8}、\q{/16}などのようなより簡潔な記法（\ac{CIDR}）を使用します。

この記法は、\ac{MSB}から始まる、マスクが持つビット数を定義するだけです。

\small
\begin{center}
\begin{tabular}{ | l | l | l | l | l | l | }
\hline
\HeaderColor マスク &
\HeaderColor ホスト &
\HeaderColor 使用可能 &
\HeaderColor ネットマスク &
\HeaderColor 16進マスク &
\HeaderColor \\
\hline
/30  & 4        & 2        & 255.255.255.252  & 0xfffffffc  & \\
\hline
/29  & 8        & 6        & 255.255.255.248  & 0xfffffff8  & \\
\hline
/28  & 16       & 14       & 255.255.255.240  & 0xfffffff0  & \\
\hline
/27  & 32       & 30       & 255.255.255.224  & 0xffffffe0  & \\
\hline
/26  & 64       & 62       & 255.255.255.192  & 0xffffffc0  & \\
\hline
/24  & 256      & 254      & 255.255.255.0    & 0xffffff00  & クラスCネットワーク \\
\hline
/23  & 512      & 510      & 255.255.254.0    & 0xfffffe00  & \\
\hline
/22  & 1024     & 1022     & 255.255.252.0    & 0xfffffc00  & \\
\hline
/21  & 2048     & 2046     & 255.255.248.0    & 0xfffff800  & \\
\hline
/20  & 4096     & 4094     & 255.255.240.0    & 0xfffff000  & \\
\hline
/19  & 8192     & 8190     & 255.255.224.0    & 0xffffe000  & \\
\hline
/18  & 16384    & 16382    & 255.255.192.0    & 0xffffc000  & \\
\hline
/17  & 32768    & 32766    & 255.255.128.0    & 0xffff8000  & \\
\hline
/16  & 65536    & 65534    & 255.255.0.0      & 0xffff0000  & クラスBネットワーク \\
\hline
/8   & 16777216 & 16777214 & 255.0.0.0        & 0xff000000  & クラスAネットワーク \\
\hline
\end{tabular}
\end{center}
\normalsize

以下は、ホストアドレスにネットワークマスクを適用してネットワークアドレスを計算する小さな例です。

\lstinputlisting[style=customc]{\CURPATH/netmask.c}

\subsection{calc\_network\_address()}

\TT{calc\_network\_address()}関数は最も単純です：
ホストアドレスとネットワークマスクをANDするだけで、ネットワークアドレスが得られます。

\lstinputlisting[caption=\Optimizing MSVC 2012 /Ob0,numbers=left,style=customasmx86]{\CURPATH/calc_network_address_MSVC_2012_Ox.asm}

22行目で最も重要な\ANDが見られます---ここでネットワークアドレスが計算されます。

\subsection{form\_IP()}

\TT{form\_IP()}関数は、すべての4バイトを32ビット値に配置するだけです。

通常の方法は以下のとおりです：

\begin{itemize}
\item 戻り値の変数を割り当てます。0に設定します。

\item 第4の（最下位の）バイトを取り、このバイトと戻り値にOR演算を適用します。
戻り値には今第4バイトが含まれています。

\item 第3バイトを取り、8ビット左にシフトします。
\TT{bb}が第3バイトの場合、\TT{0x0000bb00}のような値が得られます。
結果の値と戻り値にOR演算を適用します。
戻り値にはこれまで\TT{0x000000aa}が含まれていたので、値をORすると
\TT{0x0000bbaa}のような値が生成されます。

\item 第2バイトを取り、16ビット左にシフトします。
\TT{cc}が第2バイトの場合、\TT{0x00cc0000}のような値が得られます。
結果の値と戻り値にOR演算を適用します。
戻り値にはこれまで\TT{0x0000bbaa}が含まれていたので、値をORすると
\TT{0x00ccbbaa}のような値が生成されます。

\item 第1バイトを取り、24ビット左にシフトします。
\TT{dd}が第1バイトの場合、\TT{0xdd000000}のような値が得られます。
結果の値と戻り値にOR演算を適用します。
戻り値にはこれまで\TT{0x00ccbbaa}が含まれていたので、値をORすると
\TT{0xddccbbaa}のような値が生成されます。

\end{itemize}

そして、これが非最適化MSVC 2012によって行われる方法です：

\lstinputlisting[caption=\NonOptimizing MSVC 2012,style=customasmx86]{\CURPATH/form_IP_MSVC_2012_EN.asm}

まあ、順序は異なりますが、もちろん、演算の順序は重要ではありません。

\Optimizing MSVC 2012は本質的に同じことを行いますが、異なる方法で行います：

\lstinputlisting[caption=\Optimizing MSVC 2012 /Ob0,style=customasmx86]{\CURPATH/form_IP_MSVC_2012_Ox_EN.asm}

各バイトが戻り値の最下位8ビットに書き込まれ、その後戻り値が各ステップで1バイトずつ左にシフトされると言えるでしょう。

各入力バイトに対して4回繰り返します。

\par
以上です！残念ながら、これを行う他の方法はおそらくありません。

ビットやバイトから値を構成する命令を持つ人気のある\ac{CPU}や\ac{ISA}はありません。

通常、ビットシフトとORによって行われます。

\subsection{print\_as\_IP()}

\TT{print\_as\_IP()}は逆のことを行います：32ビット値を4バイトに分割します。

スライシングはやや単純に動作します：入力値を24、16、8、または0ビットシフトし、0番目から7番目まで（最下位バイト）のビットを取り、それで終わりです：

\lstinputlisting[caption=\NonOptimizing MSVC 2012,style=customasmx86]{\CURPATH/print_as_IP_MSVC_2012_EN.asm}

\Optimizing MSVC 2012はほぼ同じことを行いますが、入力値の不要な再ロードはありません：

\lstinputlisting[caption=\Optimizing MSVC 2012 /Ob0,style=customasmx86]{\CURPATH/print_as_IP_MSVC_2012_Ox.asm}

\subsection{form\_netmask() and set\_bit()}

\TT{form\_netmask()}は\ac{CIDR}記法からネットワークマスク値を作成します。
もちろん、ここである種の事前計算されたテーブルを使用する方がはるかに効果的でしょうが、ビットシフトを実演するために意図的にこの方法で考慮します。

また、別の関数\TT{set\_bit()}を書きます。
このような原始的な操作のために関数を作成するのはあまり良いアイデアではありませんが、すべてがどのように機能するかを理解するのは簡単でしょう。

\lstinputlisting[caption=\Optimizing MSVC 2012 /Ob0,style=customasmx86]{\CURPATH/form_netmask_MSVC_2012_Ox.asm}

\TT{set\_bit()}は原始的です：必要なビット数だけ1を左にシフトし、その後\q{input}値とORするだけです。
\TT{form\_netmask()}にはループがあります：\TT{netmask\_bits}引数で渡されたのと同じ数のビット（\ac{MSB}から開始）を設定します。

\subsection{まとめ}

以上です！
実行すると次のようになります：

\begin{lstlisting}
netmask=255.255.255.0
network address=10.1.2.0
netmask=255.0.0.0
network address=10.0.0.0
netmask=255.255.255.128
network address=10.1.2.0
netmask=255.255.255.192
network address=10.1.2.64
\end{lstlisting}